\documentclass[a4paper,fleqn]{cas-dc}

\usepackage[numbers]{natbib}
\usepackage{amsmath, amsfonts, amssymb}
\usepackage{algorithm, algpseudocode}
\usepackage{graphicx}
\usepackage{adjustbox}
\usepackage{xcolor}
\usepackage{booktabs, tabularx, multirow}
\usepackage{enumitem}
\usepackage{ragged2e}
\setlist{nosep}
\newcommand*\D{\mathop{}\!\mathrm{d}}
\definecolor{darkblue}{RGB}{0,0,139}
\newcolumntype{B}{>{\bfseries\textcolor{darkblue}}c}
\newcolumntype{C}{>{\bfseries\textcolor{darkblue}}c}

% Commenting helpers -----------------------------------------------------------
\newcommand{\sven}[1]{\textcolor{red}{[Sven: #1]}}
\newcommand{\konstantinos}[1]{\textcolor{blue}{[Konstantinos: #1]}}
\newcommand{\gijs}[1]{\textcolor{orange}{[Gijs: #1]}}
% -----------------------------------------------------------------------------
\begin{document}

\shorttitle{Pricing Swing Options in Energy Markets using DRL}
%\shortauthors{}

\title[mode=title]{Pricing Swing Options in Energy Markets using Deep Reinforcement Learning with Implicit Quantile Networks}

\author[1]{Alexander Tsoskounoglou}[]
\ead{alexander.tsoskounoglou@gmail.com}

\author[2,3]{Sven Karbach}[]
\ead{sven@karbach.org}


\author[2,3]{Konstantinos Chatziandreou}[]
\ead{k.chatziandreou@uva.nl}


%\fntext[]{The authors acknowledge support from the Amsterdam University Fund (AUF Impact Call, Fall 2024) and the AI4FinTech initiative at the University of Amsterdam.}
% Corresponding author
%\cormark[1]
%\cortext[cor1]{Corresponding author}


\affiliation[1]{%
  organization={Morgan Stanley \sven{Check once you start at MS if they accept to be mentioned }},
  addressline={},
  city={Budapest},
  postcode={},
  country={Hungary}
}

\affiliation[2]{%
  organization={Informatics Institute, University of Amsterdam},
  addressline={LAB42, Science Park 900},
  city={Amsterdam},
  postcode={1098 XH},
  country={The Netherlands}
}

\affiliation[3]{%
  organization={Korteweg-de Vries Institute for Mathematics, University of Amsterdam},
  addressline={Science Park 105–107},
  city={Amsterdam},
  postcode={1098 XG},
  country={The Netherlands}
}



% Abstract
\begin{abstract}
We present a deep reinforcement learning framework for valuing swing options in electricity markets whose spot dynamics exhibit seasonality, mean reversion, and spike risk. We formulate the swing contract as a continuous‐action stochastic control problem with volume constraints, refraction periods, and convex exercise costs. The agent observes the Markovian state of the price process and remaining flexibility, and learns optimal exercise quantities using a distributional Deep Deterministic Policy Gradient (DDPG) algorithm augmented with Implicit Quantile Networks (IQN). Our approach eliminates the need for grid-based dynamic programming and density approximations, enabling scalability to high-dimensional market models and complex contractual features. Numerical results demonstrate that the learned policies approximate optimal exercise strategies and provide accurate valuations across a wide range of market conditions. Crucially, we show that our DRL agent significantly outperforms standard Least-Squares Monte Carlo (LSM) methods in regimes with high convex costs, where the non-linear payoff structure challenges classical regression-based approximations.
\end{abstract}

\begin{keywords}
Deep Reinforcement Learning \sep Swing Options \sep Jump-Diffusion Models \sep IQN \sep Convex Costs \sep Flex Pricing
\end{keywords}


\maketitle


\section{Introduction}
\label{Introduction}

Electricity markets display distinctive features such as non-storability, strong seasonality, abrupt demand-supply imbalances and short-lived price spikes. These characteristics lead to spot price dynamics that differ fundamentally from those in traditional financial markets. A large part of the literature therefore models electricity prices with mean-reverting diffusions augmented by jump components. Early two-factor Gaussian models were introduced by~\cite{lucia2002electricity}, while~\cite{cartea2005pricing} incorporated jumps to better capture extreme movements. A major refinement was provided by~\cite{Hambly2009}, who specified the log spot price as the sum of a seasonal function, an Ornstein-Uhlenbeck process and a rapidly mean-reverting jump component. This model captures both typical price movements and spike behaviour while retaining analytical tractability through transform methods.

Swing options are central to gas and power markets because they grant the holder flexible rights to purchase energy subject to daily and cumulative constraints. Their valuation is considerably more involved than that of European-style derivatives, since the holder repeatedly decides how much to exercise at each decision time. Classical valuation approaches include tree and grid-based dynamic programming~\cite{jaillet2004valuation}, Monte-Carlo regression or dual formulations~\cite{meinshausen2004monte,carmona2004swing} and analytically tractable special cases. These approaches can be accurate but face limitations when the underlying model contains multiple stochastic factors, jump components or complex contractual features. In the jump-diffusion model from~\cite{Hambly2009}, for example, swing pricing requires a two-dimensional grid over the diffusive and jump components and the transition density of the spike process must be approximated. The computational cost increases rapidly with the dimensionality of the state space, illustrating the classical curse of dimensionality.

Reinforcement learning and neuro-dynamic programming~\cite{bertsekas1996neuro,sutton1998reinforcement} provide an attractive alternative to grid-based dynamic programming. Recent advances in financial applications include model-based and model-free hedging~\cite{halperin2017qlbs,buehler2019deephedging}, optimal stopping~\cite{becker2019deepoptimalstopping} or both~\cite{tsoskounoglou2025higherorder}. These methods are well suited to swing options because the exercise decision is naturally continuous, the horizon is long and the contractual constraints are easily incorporated through the design of the state and reward. Modern deep actor-critic algorithms with distributional critics allow for stable learning in continuous-action environments and admit efficient parallelisation.

In this paper we propose a deep reinforcement learning method for valuing swing options under jump-diffusion models. We formulate the contract as a continuous-action Markov decision process in which the agent observes the current spot price factors, the remaining volume and the time to maturity. 
We employ a distributional variant of the Deep Deterministic Policy Gradient (DDPG) algorithm, replacing the standard critic with an Implicit Quantile Network (IQN)~\cite{dabney2018implicit}. This allows the agent to learn the full distribution of returns rather than just the expectation, providing richer gradients and improved stability in the presence of heavy-tailed usage rewards driven by price spikes.
This formulation removes the need for the density approximations used in~\cite{Hambly2009} and enables scalability to high-dimensional state representations and complex contract structures including convex costs, refraction periods and operational constraints. Numerical experiments show that the learned policies provide accurate valuations and closely approximate optimal exercise patterns across a range of market conditions. In particular, we demonstrate that our method captures the non-linear value effects of convex exercise costs significantly better than standard Least-Squares Monte Carlo (LSM) baselines.

% \sven{Add more contributions here.}

Our framework therefore connects classical stochastic modelling of electricity markets with modern distributional reinforcement learning and provides a scalable method for valuing flexible energy contracts in environments characterised by volatility, spikes and intricate contractual design.

\clearpage

\section{Market and Contract Framework}

\subsection{Electricity Spot Price Dynamics: The HHK Model}

Electricity spot prices exhibit strong mean reversion, pronounced seasonality and sudden short-lived price spikes caused by supply shortages or unexpected outages. To capture these features we adopt the model of \cite{Hambly2009}, which specifies the logarithm of the spot price as the sum of a deterministic seasonal component, a diffusive mean-reverting factor and a rapidly mean-reverting jump component.

Let $S_t$ denote the spot price and define
\begin{align}\label{eq:model}
S_t = \exp\!\big(f(t) + X_t + Y_t\big),\end{align}
where $f(t)$ represents deterministic seasonality and the stochastic factors $X_t$ and $Y_t$ follow
\begin{align}
dX_t &= -\alpha X_t \, dt + \sigma\, dW_t, \label{eq:OU}\\
dY_t &= -\beta Y_t \, dt + J_t\, dN_t. \label{eq:jumpOU}
\end{align}
Here $W_t$ is a standard Brownian motion under the risk–neutral measure, $N_t$ is a Poisson process with constant intensity $\lambda$, and $J_t$ are iid jump sizes. Following~\cite{Hambly2009} and~\cite{cartea2005pricing}, we consider exponentially distributed jumps, $J_t \sim \mathrm{Exp}(1/\mu_J)$, which generate asymmetric and heavy-tailed spike behaviour. The parameter $\alpha>0$ governs the speed of mean reversion of the diffusive component \(X_t\), while \( \beta \gg \alpha \) ensures that spikes in \(Y_t\) decay quickly, consistent with empirical observations in power markets.

The seasonal function \(f(t)\) accounts for predictable periodic structure in electricity demand and supply and may be specified, for instance, as
$$
f(t) = \log(S_0) + A \cos(2\pi t),
$$
though more elaborate parametrisations can be used. Under this formulation, the pair $(X_t,Y_t)$ is Markovian and admits closed-form expressions for its conditional moments. In particular, the moment generating function of $ X_t + Y_t$ is available in semi-analytic form~\cite{Hambly2009}, enabling calibration to observed forward prices via the identity
$$
F_t^{[T]} = \mathbb{E}\left[S_T \mid \mathcal{F}_t \right],
$$
and the seasonal function $f(T)$ can be recovered to ensure consistency with market forwards, following the methodology of~\cite{cartea2005pricing} and~\cite{lucia2002electricity}.

The model~\eqref{eq:model} offers a balance between parsimony and realism: it captures the primary stylised features of electricity prices while remaining tractable. At the same time, its jump component makes classical grid-based dynamic programming significantly more demanding, particularly because accurate pricing requires approximations of the spike transition density. This motivates the reinforcement learning framework developed later in the paper, where the underlying dynamics can be simulated directly without requiring transition densities or discretisation of the state space.
\subsection{Swing Contract Specification}

Swing contracts are widely used in electricity and gas markets to provide flexible access to energy under uncertain demand and volatile spot prices. At each contractual decision time, the holder may choose how much energy to exercise, subject to both instantaneous bounds and cumulative volume restrictions. This flexibility makes swing contracts valuable but also challenging to price, as the optimal policy involves solving a multi-period stochastic control problem.

We consider a discrete set of exercise dates
\[
t_0 = 0 < t_1 < \dots < t_{T-1} < t_T, 
\qquad 
\Delta t = t_{i+1} - t_i,
\]
and denote by \(S_t\) the electricity spot price at time \(t\). At each decision time, the holder selects an exercise quantity
\[
q_t \in [q_{\min}, q_{\max}],
\]
where \(q_{\min}\) and \(q_{\max}\) represent contractual lower and upper bounds on the offtake rate. These bounds reflect operational and physical constraints such as minimum off-take requirements and maximum power delivery capacity.

\paragraph{Cumulative volume constraint.}
Over the full contract horizon, the total exercised quantity must not exceed a specified maximum:
\[
\sum_{t=0}^{T-1} q_t \le Q_{\max}.
\]
This cap ensures that the holder cannot exploit favourable price periods indefinitely and typically corresponds to an annual or seasonal consumption limit. The constraint introduces path dependence into the valuation problem, since remaining flexibility depends on past exercise decisions.

\paragraph{Operational constraints.}
In addition to quantity bounds and cumulative volume limitations, swing contracts may include further operational features. A commonly used restriction is a \emph{refraction period}, meaning that after exercising a strictly positive amount at time \(t\), the holder must wait a fixed number of periods before exercising again. More general constraints can also be incorporated, including ramping limits
\[
|q_{t+1} - q_t| \le R,
\]
monthly or daily minimum consumption levels, or limits on the number of exercise rights. These features expand the feasible set of admissible strategies and further complicate classical dynamic programming approaches, especially under jump–diffusion dynamics.

\paragraph{Payoff structure.}
Let \(K\) denote the contractual strike price. Exercising quantity \(q_t\) at time \(t\) produces the payoff
\[
\Pi_t(q_t)
=
q_t \, \max(S_t - K, 0)
 - c_{\mathrm{cost}} \, q_t^{\gamma_{\mathrm{cost}}},
\]
where \(c_{\mathrm{cost}} \ge 0\) is a cost coefficient and \(\gamma_{\mathrm{cost}} \ge 1\) is a convexity exponent. These costs capture a variety of market frictions such as transaction fees, imbalance penalties, and convex operational disutility arising from large changes in consumption. If \(\gamma_{\mathrm{cost}} = 1\), the cost is linear. If \(\gamma_{\mathrm{cost}} > 1\), the cost is convex, penalising large exercise quantities disproportionately. This non-linear structure is particularly challenging for standard regression-based valuation methods like LSM, as we demonstrate in our results.

\paragraph{Valuation objective.}
Under the risk-neutral measure, the value of the swing contract is given by
\[
V_0
=
\sup_{(q_t) \in \mathcal{Q}}
\mathbb{E}
\Bigg[
\sum_{t=0}^{T-1}
e^{-r t \Delta t}
\, \Pi_t(q_t)
\Bigg],
\]
where \(r\) denotes the risk-free interest rate, and \(\mathcal{Q}\) is the set of all exercise sequences satisfying the instantaneous bounds, cumulative volume restriction, and operational constraints. This formulation leads to a high-dimensional, continuous-control, multi-period optimisation problem that is difficult to solve using classical techniques such as tree or grid-based dynamic programming \cite{Hambly2009,jaillet2004valuation,meinshausen2004monte,carmona2004swing}. In subsequent sections we show how this structure can be naturally expressed as a Markov decision process suitable for solution by deep reinforcement learning.


\section{Benchmark: Classical Dynamic Programming}

Before introducing our reinforcement learning framework, we review the classical numerical approach for valuing swing options under the HHK model. The method, originating from \cite{Hambly2009} and closely related to the grid-based techniques of \cite{jaillet2004valuation,meinshausen2004monte,carmona2004swing}, provides an important benchmark against which we validate the performance of our deep reinforcement learning algorithm.

\subsection{State Space Discretisation and Grid Construction}

Under the HHK spot price model, the Markovian state consists of the diffusive factor \(X_t\), the spike factor \(Y_t\), the cumulative exercised volume \(Q_t\), and any auxiliary state variables modelling operational restrictions. Classical dynamic programming discretises the continuous state space by constructing structured grids:
\begin{align*}
&X_t \in \{x_1,\ldots,x_{N_X}\}, \qquad 
Y_t \in \{y_1,\ldots,y_{N_Y}\}, \qquad\\
&Q_t \in \{q_0,\ldots,q_{N_Q}\}.
\end{align*}
Even without refraction periods or ramping constraints, the two-factor HHK specification requires a two-dimensional grid for \((X,Y)\), significantly increasing complexity relative to one-factor diffusion models such as \cite{lucia2002electricity}. The grid resolution must be sufficiently fine to capture spike behaviour in the \(Y_t\) component and the mean-reverting dynamics of \(X_t\).

\subsection{Transition Probabilities Under Jump-Diffusion Dynamics}

The dynamic programming recursion requires the transition law of \((X_t,Y_t)\) over each time step \(\Delta t\). The diffusion component \(X_t\) has an exact Gaussian transition density, whereas the spike factor \(Y_t\) follows a mixture distribution arising from the Poisson jump process. As shown in \cite{Hambly2009}, the transition density of \(Y_t\) over \(\Delta t\) is approximated by truncating the Poisson distribution after \(K_{\max}\) jumps:
\[
\mathbb{P}(N_{\Delta t}=k)
= e^{-\lambda \Delta t} \frac{(\lambda \Delta t)^k}{k!}, 
\qquad k = 0,1,\ldots,K_{\max}.
\]
Conditional on the number of jumps, the distribution of \(Y_{t+\Delta t}\) admits a closed form, enabling the computation of transition probabilities to neighbouring grid points. The accuracy of this approximation depends on the grid resolution and the choice of truncation parameter \(K_{\max}\), with larger values required under high jump intensity or strong mean reversion.

\subsection{Backward Dynamic Programming Recursion}

Let \(V_t(z)\) denote the value function at time \(t\) evaluated at a grid point \(z = (x_i,y_j,q_k)\). The Bellman recursion takes the form
\[
V_t(z)
=
\max_{q \in \mathcal{A}(z)}
\Big[
\Pi_t(q)
+ 
e^{-r\Delta t}
\sum_{z'}
\mathbb{P}(Z_{t+\Delta t}=z' \mid Z_t=z, q_t=q)\,
V_{t+\Delta t}(z')
\Big],
\]
where the maximisation is taken over all admissible exercise quantities \(q\) that satisfy instantaneous and cumulative constraints. The recursion is initialised with the terminal condition \(V_T \equiv 0\) and evaluated backwards. This formulation requires evaluating transition expectations at each grid point and for each possible action, leading to computational cost proportional to
\[
O(T \cdot N_X \cdot N_Y \cdot N_Q \cdot |\mathcal{A}|).
\]

The original HHK implementation used grid sizes such as \(120 \times 60\) for \((X,Y)\), already requiring several minutes of computation for a single contract with daily exercise opportunities. As contractual complexity increases (e.g., through convex costs or operational constraints), the size of the required grid increases rapidly.

\subsection{Limitations of the Classical Grid Method}

Despite providing reliable benchmark valuations, the HHK grid method suffers from several structural limitations:
\begin{itemize}
    \item \textbf{Curse of dimensionality.} Each additional state variable introduces a new dimension in the grid. Extensions such as stochastic volatility, fuel-spread dynamics, forecast uncertainty, or renewable generation would render the grid-based method infeasible.
    \item \textbf{Jump dynamics.} Capturing the rapidly mean-reverting spike process requires fine grid resolution and careful handling of the truncated Poisson mixture, significantly increasing computational cost and potential numerical instability.
    \item \textbf{Continuous control.} Swing options allow continuous exercise quantities \(q_t \in [q_{\min},q_{\max}]\), which must be discretised in grid-based methods, introducing bias and limiting resolution.
    \item \textbf{Operational constraints.} Features such as refraction periods, ramping limits, and transaction costs increase the dimensionality of the state, exacerbating grid explosion.
\end{itemize}

These limitations underline the need for alternative approaches that avoid state-space discretisation and do not rely on transition-density approximations. In Section~4, we develop an RL framework that treats swing exercise as a continuous-control Markov decision process.

\section{Reinforcement Learning Formulation}

We now reformulate the swing option valuation problem as a continuous-control Markov decision process (MDP) and solve it using a distributional deep reinforcement learning algorithm. This approach avoids the state-space discretisation and transition-density approximations required in classical dynamic programming, while naturally accommodating continuous exercise quantities, convex cost structures and operational constraints. Our formulation follows the principles of neuro-dynamic programming \cite{bertsekas1996neuro,sutton1998reinforcement} and leverages recent advances in distributional actor--critic methods \cite{dabney2018implicit,bellemare2017distributional}.

\subsection{Markov Decision Process}

The valuation problem is represented as a discrete-time MDP on the exercise grid $\{t_0,\dots,t_{T-1}\}$. At each time $t$, the agent observes an augmented state vector
\[
s_t
=
\Big[
S_t - K, \, 
\frac{Q^{\mathrm{used}}_t}{Q_{\max}}, \,
\frac{Q^{\mathrm{ren}}_t}{Q_{\max}}, \,
\tau_t, \,
S_t, \,
X_t, \,
Y_t, \,
\delta_t, \,
\dots
\Big]
\]
where $S_t - K$ is the intrinsic value, $Q^{\mathrm{used}}_t$ is the cumulative volume exercised, $Q^{\mathrm{rem}}_t$ is the remaining volume capacity, $\tau_t$ represents the normalised time or time-to-maturity, and $(X_t, Y_t)$ are the HHK model factors. We also include $\delta_t$, the fraction of time steps since the last exercise, to encode refraction period constraints efficiently.

The agent selects a continuous exercise quantity
\[
q_t \in [q_{\min}, q_{\max}],
\]
which determines the next state through the HHK transition dynamics
\[
(X_{t+\Delta t}, Y_{t+\Delta t}) 
= 
\text{SimHHK}(X_t, Y_t; \Delta t),
\qquad
Q^{\mathrm{used}}_{t+\Delta t}
=
Q^{\mathrm{used}}_t + q_t,
\]
where $\text{SimHHK}$ denotes an Euler–Maruyama simulation of the stochastic differential equations \eqref{eq:OU}--\eqref{eq:jumpOU}. The remaining components of the state update deterministically.

The reward for exercising quantity $q_t$ at time $t$ is defined as the discounted incremental payoff,
\[
r_t(q_t)
=
e^{-r\Delta t}
\Big(
q_t \, \max(S_t - K, 0)
 - c_{\mathrm{cost}} q_t^{\gamma_{\mathrm{cost}}}
\Big).
\]
The objective of the agent is to maximise the expected cumulative reward
\[
\mathbb{E}\Bigg[
\sum_{t=0}^{T-1} r_t(q_t)
\Bigg],
\]
which coincides with the swing option value $V_0$.

\subsection{Enforcement of Swing Constraints}

Swing contracts impose restrictions on the path of admissible actions, including instantaneous bounds, cumulative volume limits and operational rules such as refraction periods. These are enforced within the environment logic:

\paragraph{Action clipping.}
At each decision time, the actor network outputs a normalised action $\tilde{q}_t \in [0,1]$, which is mapped to the admissible interval $[q_{\min},q_{\max}]$.

\paragraph{Cumulative volume bookkeeping.}
If $Q^{\mathrm{used}}_t + q_t > Q_{\max}$, the action is clipped such that $q_t = Q_{\max} - Q^{\mathrm{used}}_t$. The agent also observes the remaining capacity $Q_{\max} - Q^{\mathrm{used}}_t$ explicitly in the state vector.

\paragraph{Refraction period masking.}
If the contract prohibits exercise for $\tau_{\mathrm{ref}}$ periods following a non-zero action, the environment forces feasible actions to zero during the refraction window.

\subsection{Distributional DDPG with IQN}

To solve the MDP, we use a distributional variant of the Deep Deterministic Policy Gradient (DDPG) algorithm. We replace the standard scalar critic with an **Implicit Quantile Network (IQN)**~\cite{dabney2018implicit}. This architecture allows the agent to learn the full distribution of random returns rather than just their expectation, which is particularly beneficial in energy markets where returns are heavy-tailed due to price spikes and convex costs.

\paragraph{Deterministic actor.}
The policy is a parameterised deterministic function
\[
q_t = \pi_{\theta}(s_t),
\]
implemented by a feedforward neural network mapping states to actions in $[q_{\min},q_{\max}]$.

\paragraph{Implicit Quantile Critic.}
The critic $Z_\phi(s,a)$ models the distribution of the return $Z^\pi = \sum \gamma^t r_t$. IQN approximates this distribution by learning a quantile function $F_{Z}^{-1}(\tau)$ for $\tau \in [0,1]$. The network takes the state $s$, action $a$, and a random quantile level $\tau$ as input, and outputs the corresponding return quantile value $Q_\phi(s,a;\tau)$. 
The embedding of $\tau$ is performed using a cosine basis expansion, allowing the network to generalise across all quantile levels.

\paragraph{Training Process.}
The critic is trained to minimise the quantile regression loss (Huber loss) between the predicted quantiles and the temporal difference targets. We employ $n$-step returns to propagate rewards faster and use a replay buffer to decorrelate samples.
\[
\mathcal{L}(\phi) = \mathbb{E}_{\tau, \tau' \sim U[0,1]} \Big[ \rho_\tau^\kappa \big( r_t + \dots + \gamma^n Q_{\phi'}(s_{t+n}, \pi_{\theta'}(s_{t+n}); \tau') - Q_\phi(s_t, a_t; \tau) \big) \Big]
\]
where $\rho_\tau^\kappa$ is the quantile Huber loss.
Target networks are updated via soft averaging to ensure numerical stability.

This method combines the efficiency of deterministic policy gradients with the stability and risk-sensitivity of distributional RL, making it highly effective for the complex return landscapes of swing options.

\subsection{Implementation Details}

\paragraph{Neural network architectures.}
The actor and critic networks consist of fully connected layers with ReLU activations. The critic uses 64 cosine basis functions to embed the quantile levels.

\paragraph{State normalisation.}
State components are normalised to approximately $[-1,1]$ or $[0,1]$ to facilitate training. Specifically, price-related inputs like $S_t-K$ are scaled by the strike or historical volatility.

\paragraph{Sampling scheme.}
Trajectories are generated by simulating the HHK dynamics. We use parallel workers to collect experience into a shared replay buffer. Exploration is induced by adding noise to the actor's policy output (e.g., Ornstein-Uhlenbeck or Gaussian noise).

\section{Experimental Results}

We evaluate our DRL framework by comparing it against the standard Least-Squares Monte Carlo (LSM) method. We focus on two key aspects: pricing accuracy in standard settings and performance under convex cost structures where classical methods often struggle.

\subsection{Comparison Between LSM and DRL Valuation Methods}

We first benchmark the basic pricing accuracy using standard swing contracts without convex costs ($\gamma_{\mathrm{cost}}=1, c_{\mathrm{cost}}=0$). Both LSM and DRL are evaluated under identical HHK model parameters. LSM employs a polynomial regression basis (Chebyshev polynomials degree 7) in the state variables.
As shown in previous literature, DRL valuations closely match those produced by LSM in this linear regime, confirming that the agent successfully learns the optimal exercise boundary.

\subsection{Impact of Convex Costs}

A key advantage of our continuous-control DRL approach is its ability to handle non-linear payoff structures. We consider a swing contract where the exercise cost is given by $0.05 \cdot q_t^{\gamma_{\mathrm{cost}}}$, and we vary the convexity exponent $\gamma_{\mathrm{cost}} \in \{1.0, 1.5, 2.0, 3.0\}$.
As the cost function becomes more convex, the optimal policy shifts from a "bang-bang" (0 or Max) strategy to a nuanced interior control strategy.

Table~\ref{tab:convex_results} compares the valuations obtained by LSM and our DRL agent.

\begin{table}[h]
\centering
\caption{Comparison of Swing Option Values with Convex Costs ($c_{\mathrm{cost}}=0.05$). Values represent the mean contract price over the test set.}
\label{tab:convex_results}
\begin{tabular}{lcccc}
\toprule
$\gamma_{\mathrm{cost}}$ & Cost Type & LSM Value & DRL Value & Improvement \\
\midrule
1.0 & Linear & 2.13 & 2.13 & $\sim 0.0\%$ \\
1.5 & Convex & 1.92 & 1.97 & $+2.3\%$ \\
2.0 & Quadratic & 1.68 & 1.85 & $+10.0\%$ \\
3.0 & Cubic & 0.97 & 1.69 & \textbf{$+73.9\%$} \\
\bottomrule
\end{tabular}
\end{table}

\paragraph{Analysis.}
For linear costs ($\gamma_{\mathrm{cost}}=1.0$), both methods converge to the same value ($2.13$). This is expected as the optimal policy is simple (exercise at max capacity when profitable).
However, as convexity increases, LSM performance degrades significantly. At $\gamma_{\mathrm{cost}}=2.0$ (quadratic costs), DRL outperforms LSM by $10\%$. At $\gamma_{\mathrm{cost}}=3.0$, the gap widens to nearly $74\%$.

This discrepancy arises because LSM relies on a regression of the continuation value onto a basis of state variables. With highly convex costs, the value function becomes complex and the simple polynomial basis fails to capture the true continuation surface. Furthermore, the standard LSM optimization step at each backward node often relies on simplifications (like discretizing the action space or assuming bang-bang solutions) that are invalid for convex problems.
In contrast, the DDPG-IQN agent learns a policy network that maps states directly to continuous actions, allowing it to approximate the optimal interior solution with high precision without restrictive basis assumptions.

These results highlight the substantial value added by DRL in pricing complex, realistic energy contracts with operational frictions.

\section*{Acknowledgment}

The authors gratefully acknowledge financial support from the Amsterdam University Fund through the AUF Impact Call (Fall 2024).
Additional support from the AI4FinTech initiative at the University of Amsterdam is also gratefully acknowledged.

% References
\bibliographystyle{cas-model2-names}
\bibliography{Bibliography.bib}

\end{document}
